% notations of definitions of unit 132 that this note uses and another note of the unit states
% (made by shared/mkspec.py from shared/notation-132.json; the same TeX as in the note that states them)
\newcommand{\sieveVector}[2]{\Xi(#1;#2)}
\newcommand{\RealPlaceSizes}[4]{\operatorname{RPS}(#1,#2,#3,#4)}
\newcommand{\haarRatio}[1]{\lambda(#1)}
\newcommand{\gauge}[2]{\Gamma(#1;#2)}
\newcommand{\IsCutBessel}[7]{\operatorname{CB}(#1,#2,#3,#4,#5,#6,#7)}
\newcommand{\formConst}[1]{c_{\mathrm{f}}(#1)}
\newcommand{\gramSection}[1]{\gamma(#1)}
\newcommand{\invCorner}[1]{\operatorname{ic}(#1)}
\newcommand{\cornerDefect}[1]{\operatorname{cd}(#1)}
\newcommand{\IsBoundedInCell}[2]{#1\in\operatorname{Cell}_{\le #2}}
\newcommand{\diagLast}[1]{\operatorname{d}_{\varrho}(#1)}
\newcommand{\sectionCoeff}[2]{\mathsf{c}(#1,#2)}
\newcommand{\sectionFrame}[3]{\mathsf{f}(#1,#2,#3)}
\newcommand{\sectionMatrix}[3]{\mathsf{m}(#1,#2,#3)}
\newcommand{\headDensity}[4]{\mathsf{h}_{#1}(#4;#2,#3)}
\newcommand{\translatePairing}[5]{\langle #1(#4),#2(#3#5)\rangle}
\newcommand{\IsBessel}[2]{#1\in\operatorname{Bes}(#2)}
\newcommand{\IsCut}[1]{#1\in\operatorname{Cut}}
\newcommand{\pairingNorm}[2]{\operatorname{PN}(#1,#2)}
\newcommand{\formIntegrandOne}[6]{\mathsf{I}^{1}_{#1}(#2;#3,#4,#5)(#6)}
\newcommand{\formIntegrand}[7]{\mathsf{I}_{#1}(#2,#3;#4,#5,#6)(#7)}
\newcommand{\floorForm}[5]{\mathcal{I}_{#1}(#2,#3;#4,#5)}
\newcommand{\formModulus}[6]{\mathcal{I}^{\mathrm{mod}}_{#1}(#2,#3;#4,#5,#6)}
\newcommand{\formMajorant}[6]{\mathcal{I}^{\mathrm{maj}}_{#1}(#2,#3;#4,#5,#6)}
\newcommand{\selfNormProduct}[7]{\operatorname{SNP}_{#1}(#7;#2,#3;#4,#5,#6)}
\newcommand{\frameVolume}[1]{\Delta(#1)}
\newcommand{\frameXPart}[1]{\mathsf{x}_0(#1)}
